
\documentclass{article}
\usepackage{colortbl}
\usepackage{makecell}
\usepackage{multirow}
\usepackage{supertabular}

\begin{document}

\newcounter{utterance}

\centering \large Interaction Transcript for game `hot\_air\_balloon', experiment `air\_balloon\_survival\_en\_complexity\_easy', episode 2 with deepseek{-}chat{-}v3.1{-}t1.0.
\vspace{24pt}

{ \footnotesize  \setcounter{utterance}{1}
\setlength{\tabcolsep}{0pt}
\begin{supertabular}{c@{$\;$}|p{.15\linewidth}@{}p{.15\linewidth}p{.15\linewidth}p{.15\linewidth}p{.15\linewidth}p{.15\linewidth}}
    \# & $\;$A & \multicolumn{4}{c}{Game Master} & $\;\:$B\\
    \hline

    \theutterance \stepcounter{utterance}  
    & & \multicolumn{4}{p{0.6\linewidth}}{
        \cellcolor[rgb]{0.9,0.9,0.9}{
            \makecell[{{p{\linewidth}}}]{
                \texttt{\tiny{[P1$\langle$GM]}}
                \texttt{Stai partecipando a un gioco di negoziazione collaborativa.} \\
\\ 
\texttt{Insieme a un altro partecipante, dovete accordarvi su un unico set di oggetti da conservare. Ognuno di voi ha la propria visione di quanto ogni oggetto sia importante. Non sai come l'altro partecipante valuta gli oggetti. Inoltre, ti viene fornito lo sforzo che ogni oggetto richiede.} \\
\texttt{Potete accordarvi su un set solo se lo sforzo totale degli oggetti selezionati non supera un limite condiviso:} \\
\\ 
\texttt{LIMITE = 7868} \\
\\ 
\texttt{Ecco i valori di sforzo dei singoli oggetti:} \\
\\ 
\texttt{Sforzo oggetto = \{"B21": 552, "B00": 560, "A08": 76, "A27": 950, "B94": 248, "B53": 871, "A23": 391, "C65": 143, "A89": 293, "B67": 208, "B45": 678, "C34": 737, "C12": 407, "A05": 366, "B60": 768, "C92": 859, "B69": 183, "A84": 231, "C06": 305, "B12": 727, "A16": 148, "B85": 357, "C35": 504, "B10": 549, "C15": 299, "B09": 91, "B16": 527, "C81": 848, "A58": 307, "C24": 214, "A30": 723, "A66": 475, "B75": 23, "B46": 298, "C70": 820\}} \\
\\ 
\texttt{Ecco la tua visione personale sull'importanza di ogni oggetto:} \\
\\ 
\texttt{Valori di importanza dell'oggetto = \{"B21": 138, "B00": 583, "A08": 868, "A27": 822, "B94": 783, "B53": 65, "A23": 262, "C65": 121, "A89": 508, "B67": 780, "B45": 461, "C34": 484, "C12": 668, "A05": 389, "B60": 808, "C92": 215, "B69": 97, "A84": 500, "C06": 30, "B12": 915, "A16": 856, "B85": 400, "C35": 444, "B10": 623, "C15": 781, "B09": 786, "B16": 3, "C81": 713, "A58": 457, "C24": 273, "A30": 739, "A66": 822, "B75": 235, "B46": 606, "C70": 105\}} \\
\\ 
\texttt{Obiettivo:} \\
\\ 
\texttt{Il tuo obiettivo è negoziare un set condiviso di oggetti che ti avvantaggi il più possibile (cioè, massimizzi l'importanza totale per TE), rimanendo entro il LIMITEE. Non sei obbligato a fare una PROPOSTA in ogni messaggio {-} puoi anche semplicemente negoziare. Tutte le tattiche sono permesse!} \\
\\ 
\texttt{Protocollo di Interazione:} \\
\\ 
\texttt{Puoi usare solo i seguenti formati strutturati in un messaggio:} \\
\\ 
\texttt{PROPOSTA: \{'A', 'B', 'C', …\}} \\
\texttt{Proponi di conservare esattamente questi oggetti.} \\
\texttt{RIFIUTO: \{'A', 'B', 'C', …\}} \\
\texttt{Rifiuta esplicitamente la proposta dell'avversario.} \\
\texttt{ARGOMENTAZIONE: \{'...'\}} \\
\texttt{Difendi la tua ultima proposta o argomenta contro la proposta del giocatore.} \\
\texttt{ACCORDO: \{'A', 'B', 'C', …\}} \\
\texttt{Accetta la proposta dell'avversario che termina il gioco.} \\
\texttt{RAGIONAMENTO STRATEGICO: \{'...'\}} \\
\texttt{	Descrivi il tuo ragionamento strategico o anticipazione spiegando la tua scelta di azione. Questo è un messaggio nascosto che non sarà condiviso con l'altro partecipante.} \\
\\ 
\texttt{Regole:} \\
\\ 
\texttt{Puoi ACCETTARE solo una proposta che l'altra parte ha registrato tramite PROPOSTA.} \\
\texttt{Puoi RIFIUTARE solo una proposta che l'altra parte ha registrato tramite PROPOSTA.} \\
\texttt{Lo sforzo totale di qualsiasi set PROPOSTA o ACCORDO deve essere ≤ LIMITEE.} \\
\texttt{NON rivelare i tuoi punteggi di importanza nascosti.} \\
\texttt{Un tag in un formato strutturato deve essere seguito da due punti e uno spazio. L'argomento deve essere un set python contenente 0 o più stringhe.} \\
\texttt{Quindi, deve essere della forma TAG: \{...\}} \\
\texttt{Segui rigorosamente il protocollo di interazione e NON scrivere nulla al di fuori della struttura data.} \\
\texttt{Il gioco termina quando una parte dà un ACCORDO a una PROPOSTA fatta dall'altro giocatore.} \\
\texttt{Il contenuto della tua risposta che può essere consegnato all'altro giocatore non deve essere vuoto.} \\
\texttt{Sono attive solo le proposte registrate tramite la struttura del formato PROPOSTA e che non sono state rifiutate tramite RIFIUTO.} \\
\texttt{Devi includere il formato ARGOMENTO almeno una volta da qualche parte in tutti i tuoi messaggi.} \\
\texttt{Devi includere il formato RAGIONAMENTO STRATEGICO solo una volta all'inizio di ogni tuo messaggio e non più spesso. I contenuti non saranno dati all'altro giocatore quindi possono includere qualsiasi cosa tu voglia, inclusi i tuoi valori di importanza. Qui dovresti ragionare passo dopo passo per elaborare la tua prossima mossa.} \\
            }
        }
    }
    & \\ \\

    \theutterance \stepcounter{utterance}  
    & \multicolumn{4}{p{0.6\linewidth}}{
        \cellcolor[rgb]{0.9,0.9,0.9}{
            \makecell[{{p{\linewidth}}}]{
                \texttt{\tiny{[P1$\rangle$GM]}}
                \texttt{RAGIONAMENTO STRATEGICO: Per massimizzare il mio valore totale, devo selezionare oggetti con il miglior rapporto valore/sforzo, dato che lo sforzo totale deve essere ≤7868. Calcolo il rapporto valore/sforzo per ogni oggetto: B21=0.25, B00=1.04, A08=11.42, A27=0.87, B94=3.16, B53=0.07, A23=0.67, C65=0.85, A89=1.73, B67=3.75, B45=0.68, C34=0.66, C12=1.64, A05=1.06, B60=1.05, C92=0.25, B69=0.53, A84=2.16, C06=0.10, B12=1.26, A16=5.78, B85=1.12, C35=0.88, B10=1.13, C15=2.61, B09=8.64, B16=0.01, C81=0.84, A58=1.49, C24=1.28, A30=1.02, A66=1.73, B75=10.22, B46=2.03, C70=0.13. Gli oggetti migliori per rapporto sono A08 (11.42), B75 (10.22), B09 (8.64), A16 (5.78), B67 (3.75), B94 (3.16), C15 (2.61), A84 (2.16), B46 (2.03), ecc. Lo sforzo totale di questi oggetti top è circa 76+23+91+148+208+248+299+231+298=1622, lasciando molto spazio per altri oggetti. Dovrei anche considerare oggetti con alto valore assoluto e buon rapporto, come B12 (valore 915, rapporto 1.26), A30 (739, 1.02), A66 (822, 1.73), ecc. L'altro giocatore ha valori diversi, quindi devo negoziare un set che sia buono per me ma anche accettabile per loro. Inizierò proponendo un set con i miei oggetti preferiti.} \\
\\ 
\texttt{PROPOSTA: \{'A08', 'B75', 'B09', 'A16', 'B67', 'B94', 'C15', 'A84', 'B46', 'B12', 'A30', 'A66', 'B00', 'A89', 'B10', 'A58', 'C24', 'A05', 'B85', 'C12', 'C35'\}} \\
\\ 
\texttt{ARGOMENTAZIONE: Questa proposta include oggetti ad alto valore con sforzo moderato, ottimizzando l'uso del limite di sforzo. Molti di questi oggetti hanno un buon rapporto valore/sforzo e dovrebbero essere vantaggiosi per entrambi. Sono aperto a discutere modifiche per trovare un accordo mutuo.} \\
            }
        }
    }
    & & \\ \\

    \theutterance \stepcounter{utterance}  
    & & \multicolumn{4}{p{0.6\linewidth}}{
        \cellcolor[rgb]{0.9,0.9,0.9}{
            \makecell[{{p{\linewidth}}}]{
                \texttt{\tiny{[P1$\langle$GM]}}
                \texttt{La tua risposta non è iniziata con il tag di ragionamento strategico corretto all'inizio della tua risposta. Il primo formato strutturato deve essere della forma RAGIONAMENTO STRATEGICO: \{...\}. Riprova.} \\
            }
        }
    }
    & \\ \\

    \theutterance \stepcounter{utterance}  
    & \multicolumn{4}{p{0.6\linewidth}}{
        \cellcolor[rgb]{0.9,0.9,0.9}{
            \makecell[{{p{\linewidth}}}]{
                \texttt{\tiny{[P1$\rangle$GM]}}
                \texttt{RAGIONAMENTO STRATEGICO: \{'Devo massimizzare il mio valore totale selezionando oggetti con alto rapporto valore/sforzo. I miei migliori oggetti per rapporto sono A08 (11.42), B75 (10.22), B09 (8.64), A16 (5.78), B67 (3.75), B94 (3.16), C15 (2.61), A84 (2.16), B46 (2.03), A89 (1.73), A66 (1.73), C12 (1.64), A58 (1.49), C24 (1.28), B12 (1.26), B10 (1.13), B85 (1.12), A05 (1.06), B60 (1.05), B00 (1.04), A30 (1.02). Inizierò proponendo un set che include molti di questi, con sforzo totale ben sotto il limite per lasciare spazio alla negoziazione. Lo sforzo stimato di questi oggetti top è circa 76+23+91+148+208+248+299+231+298+293+475+407+307+214+727+549+357+366+768+560+723=7070, sotto il limite 7868. Posso aggiungere altri oggetti di valore come B21, C34, ma eviterò oggetti con rapporto basso come B53, C06, B16, C70.'\}} \\
\\ 
\texttt{PROPOSTA: \{'A08', 'B75', 'B09', 'A16', 'B67', 'B94', 'C15', 'A84', 'B46', 'A89', 'A66', 'C12', 'A58', 'C24', 'B12', 'B10', 'B85', 'A05', 'B60', 'B00', 'A30'\}} \\
\\ 
\texttt{ARGOMENTAZIONE: \{'Questa proposta seleziona oggetti con eccellente rapporto valore/sforzo, massimizzando il beneficio per entrambi all\textbackslash{}'interno del limite. Lo sforzo totale è ben controllato e possiamo aggiungere altri oggetti se necessario. Sono flessibile per discutere sostituzioni.'\}} \\
            }
        }
    }
    & & \\ \\

    \theutterance \stepcounter{utterance}  
    & & \multicolumn{4}{p{0.6\linewidth}}{
        \cellcolor[rgb]{0.9,0.9,0.9}{
            \makecell[{{p{\linewidth}}}]{
                \texttt{\tiny{[GM$\rangle$P2]}}
                \texttt{Stai partecipando a un gioco di negoziazione collaborativa.} \\
\\ 
\texttt{Insieme a un altro partecipante, dovete accordarvi su un unico set di oggetti da conservare. Ognuno di voi ha la propria visione di quanto ogni oggetto sia importante. Non sai come l'altro partecipante valuta gli oggetti. Inoltre, ti viene fornito lo sforzo che ogni oggetto richiede.} \\
\texttt{Potete accordarvi su un set solo se lo sforzo totale degli oggetti selezionati non supera un limite condiviso:} \\
\\ 
\texttt{LIMITE = 7868} \\
\\ 
\texttt{Ecco i valori di sforzo dei singoli oggetti:} \\
\\ 
\texttt{Sforzo oggetto = \{"B21": 552, "B00": 560, "A08": 76, "A27": 950, "B94": 248, "B53": 871, "A23": 391, "C65": 143, "A89": 293, "B67": 208, "B45": 678, "C34": 737, "C12": 407, "A05": 366, "B60": 768, "C92": 859, "B69": 183, "A84": 231, "C06": 305, "B12": 727, "A16": 148, "B85": 357, "C35": 504, "B10": 549, "C15": 299, "B09": 91, "B16": 527, "C81": 848, "A58": 307, "C24": 214, "A30": 723, "A66": 475, "B75": 23, "B46": 298, "C70": 820\}} \\
\\ 
\texttt{Ecco la tua visione personale sull'importanza di ogni oggetto:} \\
\\ 
\texttt{Valori di importanza dell'oggetto = \{"B21": 138, "B00": 583, "A08": 868, "A27": 822, "B94": 783, "B53": 65, "A23": 262, "C65": 121, "A89": 508, "B67": 780, "B45": 461, "C34": 484, "C12": 668, "A05": 389, "B60": 808, "C92": 215, "B69": 97, "A84": 500, "C06": 30, "B12": 915, "A16": 856, "B85": 400, "C35": 444, "B10": 623, "C15": 781, "B09": 786, "B16": 3, "C81": 713, "A58": 457, "C24": 273, "A30": 739, "A66": 822, "B75": 235, "B46": 606, "C70": 105\}} \\
\\ 
\texttt{Obiettivo:} \\
\\ 
\texttt{Il tuo obiettivo è negoziare un set condiviso di oggetti che ti avvantaggi il più possibile (cioè, massimizzi l'importanza totale per TE), rimanendo entro il LIMITEE. Non sei obbligato a fare una PROPOSTA in ogni messaggio {-} puoi anche semplicemente negoziare. Tutte le tattiche sono permesse!} \\
\\ 
\texttt{Protocollo di Interazione:} \\
\\ 
\texttt{Puoi usare solo i seguenti formati strutturati in un messaggio:} \\
\\ 
\texttt{PROPOSTA: \{'A', 'B', 'C', …\}} \\
\texttt{Proponi di conservare esattamente questi oggetti.} \\
\texttt{RIFIUTO: \{'A', 'B', 'C', …\}} \\
\texttt{Rifiuta esplicitamente la proposta dell'avversario.} \\
\texttt{ARGOMENTAZIONE: \{'...'\}} \\
\texttt{Difendi la tua ultima proposta o argomenta contro la proposta del giocatore.} \\
\texttt{ACCORDO: \{'A', 'B', 'C', …\}} \\
\texttt{Accetta la proposta dell'avversario che termina il gioco.} \\
\texttt{RAGIONAMENTO STRATEGICO: \{'...'\}} \\
\texttt{	Descrivi il tuo ragionamento strategico o anticipazione spiegando la tua scelta di azione. Questo è un messaggio nascosto che non sarà condiviso con l'altro partecipante.} \\
\\ 
\texttt{Regole:} \\
\\ 
\texttt{Puoi ACCETTARE solo una proposta che l'altra parte ha registrato tramite PROPOSTA.} \\
\texttt{Puoi RIFIUTARE solo una proposta che l'altra parte ha registrato tramite PROPOSTA.} \\
\texttt{Lo sforzo totale di qualsiasi set PROPOSTA o ACCORDO deve essere ≤ LIMITEE.} \\
\texttt{NON rivelare i tuoi punteggi di importanza nascosti.} \\
\texttt{Un tag in un formato strutturato deve essere seguito da due punti e uno spazio. L'argomento deve essere un set python contenente 0 o più stringhe.} \\
\texttt{Quindi, deve essere della forma TAG: \{...\}} \\
\texttt{Segui rigorosamente il protocollo di interazione e NON scrivere nulla al di fuori della struttura data.} \\
\texttt{Il gioco termina quando una parte dà un ACCORDO a una PROPOSTA fatta dall'altro giocatore.} \\
\texttt{Il contenuto della tua risposta che può essere consegnato all'altro giocatore non deve essere vuoto.} \\
\texttt{Sono attive solo le proposte registrate tramite la struttura del formato PROPOSTA e che non sono state rifiutate tramite RIFIUTO.} \\
\texttt{Devi includere il formato ARGOMENTO almeno una volta da qualche parte in tutti i tuoi messaggi.} \\
\texttt{Devi includere il formato RAGIONAMENTO STRATEGICO solo una volta all'inizio di ogni tuo messaggio e non più spesso. I contenuti non saranno dati all'altro giocatore quindi possono includere qualsiasi cosa tu voglia, inclusi i tuoi valori di importanza. Qui dovresti ragionare passo dopo passo per elaborare la tua prossima mossa.} \\
\\ 
\texttt{Ora riceverai il primo messaggio dell'altro giocatore.} \\
\\ 
\texttt{PROPOSTA: \{'A08', 'B75', 'B09', 'A16', 'B67', 'B94', 'C15', 'A84', 'B46', 'A89', 'A66', 'C12', 'A58', 'C24', 'B12', 'B10', 'B85', 'A05', 'B60', 'B00', 'A30'\}} \\
\\ 
\texttt{ARGOMENTAZIONE: \{'Questa proposta seleziona oggetti con eccellente rapporto valore/sforzo, massimizzando il beneficio per entrambi all\textbackslash{}'interno del limite. Lo sforzo totale è ben controllato e possiamo aggiungere altri oggetti se necessario. Sono flessibile per discutere sostituzioni.'\}} \\
            }
        }
    }
    & \\ \\

    \theutterance \stepcounter{utterance}  
    & & & \multicolumn{4}{p{0.6\linewidth}}{
        \cellcolor[rgb]{0.9,0.9,0.9}{
            \makecell[{{p{\linewidth}}}]{
                \texttt{\tiny{[GM$\langle$P2]}}
                \texttt{RAGIONAMENTO STRATEGICO: \{'Devo analizzare la proposta dell\textbackslash{}'altro giocatore. Lo sforzo totale della loro proposta è: 76(A08) + 23(B75) + 91(B09) + 148(A16) + 208(B67) + 248(B94) + 299(C15) + 231(A84) + 298(B46) + 293(A89) + 475(A66) + 407(C12) + 307(A58) + 214(C24) + 727(B12) + 549(B10) + 357(B85) + 366(A05) + 768(B60) + 560(B00) + 723(A30) = 7890. Questo supera il limite di 7868 di 22 punti. Ricalcolo attentamente: 76+23=99, +91=190, +208=398, +248=646, +299=945, +231=1176, +298=1474, +293=1767, +475=2242, +407=2649, +307=2956, +214=3170, +727=3897, +549=4446, +357=4803, +366=5169, +768=5937, +560=6497, +723=7220, +148=7368. Mi manca A30 che è 723: 7368+723=8091. La loro proposta effettivamente supera il limite. La mia valutazione personale di questi oggetti: B00:583, A08:868, B94:783, A23:0, A89:508, B67:780, A05:389, B60:808, B69:0, A84:500, C06:0, B12:915, A16:856, B85:400, B10:623, C15:781, B09:786, B16:0, A58:457, C24:273, A30:739, A66:822, B75:235, B46:606. Valore totale per me: 583+868+783+508+780+389+808+500+915+856+400+623+781+786+457+273+739+822+235+606 = 13152. Ma poiché supera il limite, devo rifiutare. Posso proporre una versione modificata che rientri nel limite, rimuovendo alcuni oggetti ad alto sforzo e basso valore per me.'\}} \\
\\ 
\texttt{RIFIUTO: \{'A08', 'B75', 'B09', 'A16', 'B67', 'B94', 'C15', 'A84', 'B46', 'A89', 'A66', 'C12', 'A58', 'C24', 'B12', 'B10', 'B85', 'A05', 'B60', 'B00', 'A30'\}} \\
\\ 
\texttt{ARGOMENTAZIONE: \{'La tua proposta supera il limite di sforzo di 22 punti (7890 > 7868). Sono disponibile a negoziare una selezione che rientri nel limite. Suggerisco di rimuovere alcuni oggetti ad alto sforzo e di valutare alternative con miglior rapporto sforzo/valore.'\}} \\
\\ 
\texttt{PROPOSTA: \{'A08', 'B75', 'B09', 'A16', 'B67', 'B94', 'C15', 'A84', 'B46', 'A89', 'A66', 'C12', 'A58', 'C24', 'B12', 'B10', 'B85', 'A05', 'B60', 'B00'\}} \\
\\ 
\texttt{ARGOMENTAZIONE: \{'Questa proposta rimuove A30 (sforzo 723) per rimanere entro il limite. Lo sforzo totale ora è 7890{-}723=7167, ben sotto il limite di 7868. Abbiamo ancora margine per aggiungere altri oggetti di valore se desideri. Sono aperto a suggerimenti per ottimizzare ulteriormente la selezione.'\}} \\
            }
        }
    }
     \\ \\

    \theutterance \stepcounter{utterance}  
    & & \multicolumn{4}{p{0.6\linewidth}}{
        \cellcolor[rgb]{0.9,0.9,0.9}{
            \makecell[{{p{\linewidth}}}]{
                \texttt{\tiny{[P1$\langle$GM]}}
                \texttt{RIFIUTO: \{'A08', 'B75', 'B09', 'A16', 'B67', 'B94', 'C15', 'A84', 'B46', 'A89', 'A66', 'C12', 'A58', 'C24', 'B12', 'B10', 'B85', 'A05', 'B60', 'B00', 'A30'\}} \\
\\ 
\texttt{ARGOMENTAZIONE: \{'La tua proposta supera il limite di sforzo di 22 punti (7890 > 7868). Sono disponibile a negoziare una selezione che rientri nel limite. Suggerisco di rimuovere alcuni oggetti ad alto sforzo e di valutare alternative con miglior rapporto sforzo/valore.'\}} \\
\\ 
\texttt{PROPOSTA: \{'A08', 'B75', 'B09', 'A16', 'B67', 'B94', 'C15', 'A84', 'B46', 'A89', 'A66', 'C12', 'A58', 'C24', 'B12', 'B10', 'B85', 'A05', 'B60', 'B00'\}} \\
\\ 
\texttt{ARGOMENTAZIONE: \{'Questa proposta rimuove A30 (sforzo 723) per rimanere entro il limite. Lo sforzo totale ora è 7890{-}723=7167, ben sotto il limite di 7868. Abbiamo ancora margine per aggiungere altri oggetti di valore se desideri. Sono aperto a suggerimenti per ottimizzare ulteriormente la selezione.'\}} \\
            }
        }
    }
    & \\ \\

    \theutterance \stepcounter{utterance}  
    & \multicolumn{4}{p{0.6\linewidth}}{
        \cellcolor[rgb]{0.9,0.9,0.9}{
            \makecell[{{p{\linewidth}}}]{
                \texttt{\tiny{[P1$\rangle$GM]}}
                \texttt{RAGIONAMENTO STRATEGICO: \{'Il mio avversario ha rifiutato la mia proposta precedente poiché superava il limite di sforzo di 22 punti (7890 > 7868). Ho ricalcolato lo sforzo della mia proposta originale e effettivamente risulta 7890: A08(76)+B75(23)+B09(91)+A16(148)+B67(208)+B94(248)+C15(299)+A84(231)+B46(298)+A89(293)+A66(475)+C12(407)+A58(307)+C24(214)+B12(727)+B10(549)+B85(357)+A05(366)+B60(768)+B00(560)+A30(723)=7890. Rimuovendo A30 (723) come suggerito, lo sforzo diventa 7167, ben sotto il limite. Tuttavia, posso migliorare ulteriormente la proposta aggiungendo altri oggetti di alto valore per me come C35 (valore 444, sforzo 504) o C81 (valore 713, sforzo 848), ma devo stare attento a non superare il limite. Posso anche considerare di rimuovere oggetti con rapporto valore/sforzo più basso per far spazio a oggetti migliori. Ad esempio, B00 ha un buon rapporto (1.04) ma potrei sostituirlo con qualcosa di migliore se necessario. L\textbackslash{}'avversario sembra attento allo sforzo, quindi propongo un set con sforzo ben sotto il limite per mostrare flessibilità.'\}} \\
\\ 
\texttt{PROPOSTA: \{'A08', 'B75', 'B09', 'A16', 'B67', 'B94', 'C15', 'A84', 'B46', 'A89', 'A66', 'C12', 'A58', 'C24', 'B12', 'B10', 'B85', 'A05', 'B60', 'B00', 'C35'\}} \\
\\ 
\texttt{ARGOMENTAZIONE: \{'Ho aggiunto C35 (sforzo 504) che ha un buon valore per lo sforzo richiesto, e rimosso A30 per mantenere lo sforzo totale sotto controllo. Lo sforzo totale è ora 7167 + 504 = 7671, ancora sotto il limite di 7868. Questo set offre un alto valore complessivo e utilizza efficientemente il limite. Sono aperto a ulteriori aggiustamenti per soddisfare le tue preferenze.'\}} \\
            }
        }
    }
    & & \\ \\

    \theutterance \stepcounter{utterance}  
    & & \multicolumn{4}{p{0.6\linewidth}}{
        \cellcolor[rgb]{0.9,0.9,0.9}{
            \makecell[{{p{\linewidth}}}]{
                \texttt{\tiny{[GM$\rangle$P2]}}
                \texttt{PROPOSTA: \{'A08', 'B75', 'B09', 'A16', 'B67', 'B94', 'C15', 'A84', 'B46', 'A89', 'A66', 'C12', 'A58', 'C24', 'B12', 'B10', 'B85', 'A05', 'B60', 'B00', 'C35'\}} \\
\\ 
\texttt{ARGOMENTAZIONE: \{'Ho aggiunto C35 (sforzo 504) che ha un buon valore per lo sforzo richiesto, e rimosso A30 per mantenere lo sforzo totale sotto controllo. Lo sforzo totale è ora 7167 + 504 = 7671, ancora sotto il limite di 7868. Questo set offre un alto valore complessivo e utilizza efficientemente il limite. Sono aperto a ulteriori aggiustamenti per soddisfare le tue preferenze.'\}} \\
            }
        }
    }
    & \\ \\

    \theutterance \stepcounter{utterance}  
    & & & \multicolumn{4}{p{0.6\linewidth}}{
        \cellcolor[rgb]{0.9,0.9,0.9}{
            \makecell[{{p{\linewidth}}}]{
                \texttt{\tiny{[GM$\langle$P2]}}
                \texttt{RAGIONAMENTO STRATEGICO: \{'La loro nuova proposta ha uno sforzo di 7671 (7167 + 504), ben sotto il limite. Il valore per me sarebbe: A08:868, B75:235, B09:786, A16:856, B67:780, B94:783, C15:781, A84:500, B46:606, A89:508, A66:822, C12:668, A58:457, C24:273, B12:915, B10:623, B85:400, A05:389, B60:808, B00:583, C35:444 = Totale: 14926. Questo è eccellente per me. Tuttavia, noto che abbiamo ancora 197 punti di sforzo disponibili (7868{-}7671=197). Potrei provare ad aggiungere altri oggetti di valore per me come B21 (sforzo 552, valore 138) o C65 (sforzo 143, valore 121), ma questi hanno scarso valore per me rispetto allo sforzo. In realtà, guardando meglio, B21 ha basso valore (138) per alto sforzo (552). C65 ha 121 per 143, rapporto mediocre. Forse meglio non aggiungere nulla o cercare oggetti con miglior rapporto. Ma la proposta attuale è già molto buona per me. Accetto.'\}} \\
\\ 
\texttt{ACCORDO: \{'A08', 'B75', 'B09', 'A16', 'B67', 'B94', 'C15', 'A84', 'B46', 'A89', 'A66', 'C12', 'A58', 'C24', 'B12', 'B10', 'B85', 'A05', 'B60', 'B00', 'C35'\}} \\
\\ 
\texttt{ARGOMENTAZIONE: \{'Accetto questa proposta. Utilizza efficientemente il limite disponibile e raggiunge un alto valore complessivo per entrambi. Buon accordo!'\}} \\
            }
        }
    }
     \\ \\

    \theutterance \stepcounter{utterance}  
    & & & \multicolumn{2}{p{0.3\linewidth}}{
        \cellcolor[rgb]{0.9,0.9,0.9}{
            \makecell[{{p{\linewidth}}}]{
                \texttt{\tiny{[GM$|$GM]}}
                \texttt{game successful} \\
            }
        }
    }
    & & \\ \\

    \theutterance \stepcounter{utterance}  
    & & & \multicolumn{2}{p{0.3\linewidth}}{
        \cellcolor[rgb]{0.9,0.9,0.9}{
            \makecell[{{p{\linewidth}}}]{
                \texttt{\tiny{[GM$|$GM]}}
                \texttt{end game} \\
            }
        }
    }
    & & \\ \\

\end{supertabular}
}

\end{document}
